% answers_short.tex —— 2001 年·上海卷·文科 · 简短答案速查
% 由 tools/make_short_answers.py 自 answers.tex 的 \gksolhead/\gkanswer 提取，
% 并与 2001/2001_数学真题_上海文_答案要点.md 逐题比对后生成，勿手改（改源 MD 后重跑）。

\begin{gknotice}
\textbf{使用说明}：本节只给最终答案便于快速核对，完整推导见解析版。答案为编写者独立推导，\textbf{非官方标准答案}。
\end{gknotice}

\gkpart{填空题}
\begin{itemize}[leftmargin=2.6em]
  \item[\textbf{1.}] $3$
  \item[\textbf{2.}] $153$
  \item[\textbf{3.}] $x^2-4y^2=1$
  \item[\textbf{4.}] $1$
  \item[\textbf{5.}] $\left(0,\dfrac14\right)$
  \item[\textbf{6.}] $(0,7)$
  \item[\textbf{7.}] $7$
  \item[\textbf{8.}] $0$
  \item[\textbf{9.}] $\left[0,\dfrac{2\pi}{3}\right]$
  \item[\textbf{10.}] $A_3$
  \item[\textbf{11.}] 两圆方程相减得到的方程表示两圆的根轴，即两圆上幂相等的点的轨迹；相交时为公共弦所在直线。
  \item[\textbf{12.}] 按题面数据在右图网格中绘制对应的年平均土地沙化面积图，见矢量图源。
\end{itemize}

\gkpart{选择题}
\begin{center}\begin{tabular}{ccccc}
\toprule
题号 & 13 & 14 & 15 & 16 \\
\midrule
答案 & C & A & D & D \\
\bottomrule
\end{tabular}\end{center}

\gkpart{解答题}
\begin{itemize}[leftmargin=2.6em]
  \item[\textbf{17.}] $c=\sqrt{21}$ 或 $c=\sqrt{61}$。
  \item[\textbf{18.}] $2$ 或 $\dfrac72$。
  \item[\textbf{19.}] (1) 见解析；(2) 二面角大小为 $\arctan(2\sqrt2)$
  \item[\textbf{20.}] (1) $M_a=\{1\}$，$P=0$；(2) 例如 $z=e^{i\pi/3}$
  \item[\textbf{21.}] (1) $f(0)=1$，表示不使用水清洗时，清洗后残留量与清洗前相等；(2) 可取定义域 $[0,+\infty)$；$f(0)=1$，$0<f(x)\le1$，且 $f$ 随 $x$ 增大而减小，通常还应有 $\displaystyle\lim_{x\to+\infty}f(x)=0$；(3) $0<a<2\sqrt2$ 时一次清洗较少；$a=2\sqrt2$ 时两种方案相同；$a>2\sqrt2$ 时分两次清洗较少
  \item[\textbf{22.}] (1) $\{x_n\}=\left\{\dfrac{49}{65},\dfrac{11}{19},\dfrac15,-1\right\}$，输出 $-1$ 后停止；(2) $x_0=1$ 或 $x_0=2$；(3) 不存在这样的 $x_0$
\end{itemize}
